\documentclass{beamer}

\usepackage[T1]{fontenc}
\usepackage[utf8]{inputenc}
\usepackage{lmodern}
\usepackage{amsmath,amssymb,amsfonts}
\usepackage{physics}
\usepackage{bm}
\usepackage{mathtools}
\usepackage{quantikz}

\title{The Quantum Fourier Transform}
\subtitle{Mathematics, Circuit Decomposition, and Implementation}
\author{}
\date{}

\begin{document}

%------------------------------------------------
\frame{\titlepage}

%------------------------------------------------
\begin{frame}{Outline}
\begin{enumerate}
\item Motivation and definition
\item From the classical DFT to the QFT
\item Derivation of the product form
\item Circuit decomposition with Hadamards and controlled phases
\item Worked 2- and 3-qubit examples
\item Approximate QFT and complexity
\item Connection to quantum phase estimation
\item Implementation logic and simulation remarks
\end{enumerate}
\end{frame}

%------------------------------------------------
\begin{frame}{Motivation}
The Quantum Fourier Transform (QFT) is one of the central primitives in quantum computing.

It appears in:
\begin{itemize}
\item Shor's factoring algorithm,
\item order finding,
\item period finding,
\item quantum phase estimation,
\item several quantum simulation algorithms.
\end{itemize}

The QFT is the quantum analogue of the discrete Fourier transform, but it acts on the amplitudes of a quantum state.
\end{frame}

%------------------------------------------------
\begin{frame}{Classical Discrete Fourier Transform}
Given a vector \(x=(x_0,x_1,\dots,x_{N-1})\), the discrete Fourier transform is
\[
y_k = \frac{1}{\sqrt N}\sum_{j=0}^{N-1} x_j e^{2\pi i jk/N},
\qquad k=0,\dots,N-1.
\]

Equivalently, the DFT matrix is
\[
F_N = \frac{1}{\sqrt N}
\begin{pmatrix}
1 & 1 & 1 & \cdots & 1 \\
1 & \omega & \omega^2 & \cdots & \omega^{N-1} \\
1 & \omega^2 & \omega^4 & \cdots & \omega^{2(N-1)} \\
\vdots & \vdots & \vdots & \ddots & \vdots \\
1 & \omega^{N-1} & \omega^{2(N-1)} & \cdots & \omega^{(N-1)^2}
\end{pmatrix}
\]
with
\[
\omega = e^{2\pi i/N}.
\]
\end{frame}

%------------------------------------------------
\begin{frame}{Definition of the QFT}
For \(N=2^n\), the QFT is the unitary transformation
\[
\operatorname{QFT}_N \ket{x}
=
\frac{1}{\sqrt N}
\sum_{k=0}^{N-1}
e^{2\pi i xk/N}\ket{k},
\qquad x=0,\dots,N-1.
\]

Thus the QFT maps a computational basis state \(\ket{x}\) to a coherent superposition of all basis states with phases determined by \(x\).

Because the map is linear, for a general state
\[
\ket{\psi}=\sum_{x=0}^{N-1} c_x \ket{x},
\]
we have
\[
\operatorname{QFT}_N \ket{\psi}
=
\sum_{x=0}^{N-1} c_x \operatorname{QFT}_N\ket{x}.
\]
\end{frame}

%------------------------------------------------
\begin{frame}{Binary Representation of the Input}
Let
\[
x = x_1 x_2 \dots x_n
\]
be the binary representation of the integer \(x\), meaning
\[
x = \sum_{j=1}^n x_j 2^{n-j},
\qquad x_j\in\{0,1\}.
\]

We also write binary fractions as
\[
0.x_j x_{j+1}\dots x_n
=
\frac{x_j}{2}+\frac{x_{j+1}}{2^2}+\cdots+\frac{x_n}{2^{n-j+1}}.
\]

This notation is crucial because the QFT phases factor naturally into such binary fractions.
\end{frame}

%------------------------------------------------
\begin{frame}{Step-by-Step Derivation of the Product Form I}
Start from
\[
\operatorname{QFT}_N \ket{x}
=
\frac{1}{\sqrt{2^n}}
\sum_{k=0}^{2^n-1}
e^{2\pi i xk/2^n}\ket{k}.
\]

Write the output label \(k\) in binary:
\[
k = k_1 k_2 \dots k_n,
\qquad
k = \sum_{m=1}^n k_m 2^{n-m}.
\]

Then
\[
\frac{xk}{2^n}
=
x \sum_{m=1}^n \frac{k_m}{2^m}.
\]

Hence
\[
e^{2\pi i xk/2^n}
=
\prod_{m=1}^n e^{2\pi i x k_m / 2^m}.
\]
\end{frame}

%------------------------------------------------
\begin{frame}{Step-by-Step Derivation of the Product Form II}
Because each \(k_m\in\{0,1\}\), the sum over all \(k\) factorizes:
\[
\operatorname{QFT}_N \ket{x}
=
\frac{1}{2^{n/2}}
\bigotimes_{m=1}^n
\left(
\ket{0} + e^{2\pi i x/2^m}\ket{1}
\right).
\]

Now only the fractional part of \(x/2^m\) matters in the phase, and that fractional part is exactly
\[
0.x_{n-m+1}x_{n-m+2}\dots x_n.
\]

Thus
\[
\operatorname{QFT}_N \ket{x_1 x_2 \dots x_n}
=
\frac{1}{2^{n/2}}
\left(\ket{0}+e^{2\pi i\,0.x_n}\ket{1}\right)
\otimes
\left(\ket{0}+e^{2\pi i\,0.x_{n-1}x_n}\ket{1}\right)
\otimes \cdots
\]
\[
\cdots \otimes
\left(\ket{0}+e^{2\pi i\,0.x_1x_2\dots x_n}\ket{1}\right).
\]

The output order is therefore reversed relative to the input qubit order.
\end{frame}

%------------------------------------------------
\begin{frame}{Final Product Form}
The standard product form is
\[
\operatorname{QFT}_N \ket{x_1x_2\dots x_n}
=
\frac{1}{2^{n/2}}
\left(\ket{0}+e^{2\pi i\,0.x_n}\ket{1}\right)
\otimes
\left(\ket{0}+e^{2\pi i\,0.x_{n-1}x_n}\ket{1}\right)
\otimes \cdots
\]
\[
\cdots \otimes
\left(\ket{0}+e^{2\pi i\,0.x_1x_2\dots x_n}\ket{1}\right).
\]

This factorization is the key reason the QFT can be implemented efficiently.

It shows that the QFT can be built from:
\begin{itemize}
\item Hadamard gates,
\item controlled phase rotations,
\item a final qubit reversal using SWAP gates.
\end{itemize}
\end{frame}

%------------------------------------------------
\begin{frame}{Basic Gates Used in the QFT}
The main one-qubit gate is the Hadamard
\[
H = \frac{1}{\sqrt 2}
\begin{pmatrix}
1 & 1 \\
1 & -1
\end{pmatrix},
\]
which creates equal superpositions.

The phase rotations are
\[
R_k =
\begin{pmatrix}
1 & 0 \\
0 & e^{2\pi i / 2^k}
\end{pmatrix}.
\]

Controlled-\(R_k\) gates apply \(R_k\) to a target qubit only if the control qubit is \(\ket{1}\).
\end{frame}

%------------------------------------------------
\begin{frame}{How the Circuit Arises}
The first output qubit in the product form contains the phase
\[
e^{2\pi i\,0.x_1x_2\dots x_n}.
\]

This is produced by:
\begin{itemize}
\item a Hadamard on the first qubit, generating \(\frac{1}{\sqrt2}(\ket0+\ket1)\),
\item then a sequence of controlled phase gates from the later qubits, adding the binary-fraction phase.
\end{itemize}

The same pattern repeats recursively for the remaining qubits.

Thus the QFT circuit is naturally built qubit by qubit.
\end{frame}

%------------------------------------------------
\begin{frame}{Two-Qubit QFT: Mathematical Form}
For \(n=2\), we have \(N=4\), and
\[
\operatorname{QFT}_4 \ket{x_1x_2}
=
\frac{1}{2}
\left(\ket0 + e^{2\pi i\,0.x_2}\ket1\right)
\otimes
\left(\ket0 + e^{2\pi i\,0.x_1x_2}\ket1\right).
\]

This requires:
\begin{itemize}
\item Hadamard on the first qubit,
\item controlled-\(R_2\),
\item Hadamard on the second qubit,
\item then a SWAP to correct the reversed output order.
\end{itemize}
\end{frame}

%------------------------------------------------
\begin{frame}{Two-Qubit QFT Circuit}
\centering
\begin{quantikz}
\lstick{$q_0$} & \gate{H} & \ctrl{1} & \qw & \swap{1} & \qw \\
\lstick{$q_1$} & \qw & \gate{R_2} & \gate{H} & \targX{}{} & \qw
\end{quantikz}


Here:
\begin{itemize}
\item \(H\) creates the superposition,
\item \(R_2=\mathrm{diag}(1,e^{i\pi/2})\),
\item the SWAP reverses the qubit order.
\end{itemize}
\end{frame}

%------------------------------------------------
\begin{frame}{Three-Qubit QFT: Mathematical Structure}
For \(n=3\),
\[
\operatorname{QFT}_8 \ket{x_1x_2x_3}
=
\frac{1}{\sqrt 8}
\left(\ket0+e^{2\pi i\,0.x_3}\ket1\right)
\otimes
\left(\ket0+e^{2\pi i\,0.x_2x_3}\ket1\right)
\otimes
\left(\ket0+e^{2\pi i\,0.x_1x_2x_3}\ket1\right).
\]

The first qubit therefore requires phase contributions corresponding to
\[
\frac{x_2}{2} + \frac{x_3}{4},
\]
and this is exactly what the controlled \(R_2\) and \(R_3\) gates generate.
\end{frame}

%------------------------------------------------
\begin{frame}{Three-Qubit QFT Circuit}
\centering
\begin{quantikz}
\lstick{$q_0$} & \gate{H} & \ctrl{1}    & \ctrl{2}    & \qw      & \qw        & \qw      & \swap{2}   & \qw \\
\lstick{$q_1$} & \qw      & \gate{R_2}  & \qw         & \gate{H} & \ctrl{1}   & \qw      & \qw        & \qw \\
\lstick{$q_2$} & \qw      & \qw         & \gate{R_3}  & \qw      & \gate{R_2} & \gate{H} & \targX{}{} & \qw
\end{quantikz}

\vspace{0.5em}

Conceptually:
\begin{enumerate}
\item \(H\) on qubit 0, then controlled \(R_2,R_3\)
\item \(H\) on qubit 1, then controlled \(R_2\)
\item \(H\) on qubit 2
\item final reversal of qubit order
\end{enumerate}
\end{frame}

%------------------------------------------------
\begin{frame}{General \(n\)-Qubit QFT Circuit}
For qubit \(j\), counting from the top:

\begin{enumerate}
\item Apply a Hadamard \(H\)
\item For each later qubit \(k>j\), apply a controlled \(R_{k-j+1}\)
\item Continue recursively
\item At the end, reverse qubit order with SWAP gates
\end{enumerate}

So the general structure is:
\[
H \to \text{controlled phases} \to H \to \text{controlled phases} \to \cdots \to \text{SWAP network}.
\]
\end{frame}

%------------------------------------------------
\begin{frame}{General Pattern in Quantikz}
\centering
\begin{quantikz}
\lstick{$q_0$} & \gate{H} & \ctrl{1} & \ctrl{2} & \ctrl{3} & \qw \\
\lstick{$q_1$} & \qw & \gate{R_2} & \ctrl{1} & \ctrl{2} & \qw \\
\lstick{$q_2$} & \qw & \qw & \gate{R_2} & \ctrl{1} & \qw \\
\lstick{$q_3$} & \qw & \qw & \qw & \gate{R_2} & \gate{H}
\end{quantikz}

\vspace{0.5em}

This is a schematic pattern rather than a fully expanded full-\(n\) circuit. The key point is the nested structure of Hadamards and controlled phase gates.
\end{frame}

%------------------------------------------------
\begin{frame}{Worked 3-Qubit Phase Example I}
Take the computational basis input
\[
\ket{x}=\ket{101}.
\]

This means
\[
x_1=1,\qquad x_2=0,\qquad x_3=1,
\qquad x=5.
\]

The product form gives
\[
\operatorname{QFT}_8 \ket{101}
=
\frac{1}{\sqrt8}
\left(\ket0 + e^{2\pi i\,0.1}\ket1\right)
\otimes
\left(\ket0 + e^{2\pi i\,0.01}\ket1\right)
\otimes
\left(\ket0 + e^{2\pi i\,0.101}\ket1\right).
\]
\end{frame}

%------------------------------------------------
\begin{frame}{Worked 3-Qubit Phase Example II}
Now evaluate the binary fractions:
\[
0.1 = \frac12,
\qquad
0.01 = \frac14,
\qquad
0.101 = \frac12 + \frac18 = \frac58.
\]

Hence
\[
e^{2\pi i\,0.1} = e^{i\pi} = -1,
\]
\[
e^{2\pi i\,0.01} = e^{i\pi/2} = i,
\]
\[
e^{2\pi i\,0.101} = e^{2\pi i(5/8)} = e^{5\pi i/4}.
\]

Therefore
\[
\operatorname{QFT}_8 \ket{101}
=
\frac{1}{\sqrt8}
(\ket0-\ket1)\otimes(\ket0+i\ket1)\otimes(\ket0+e^{5\pi i/4}\ket1).
\]
\end{frame}

%------------------------------------------------
\begin{frame}{Worked 3-Qubit Phase Example III}
This example shows exactly how the QFT encodes the binary digits of the input into relative phases.

This phase encoding is what makes the QFT so useful in:
\begin{itemize}
\item phase estimation,
\item period finding,
\item extracting eigenphases of unitaries.
\end{itemize}

In other words, the QFT converts arithmetic structure into measurable interference structure.
\end{frame}

%------------------------------------------------
\begin{frame}{Inverse QFT}
The inverse QFT is obtained by reversing the circuit and taking Hermitian conjugates:
\begin{itemize}
\item replace each \(R_k\) by \(R_k^\dagger\),
\item reverse the order of the gates,
\item keep the same SWAP network.
\end{itemize}

Thus the inverse QFT is also efficient.

In practice, many algorithms use the inverse QFT rather than the forward QFT.
\end{frame}

%------------------------------------------------
\begin{frame}{Complexity of the QFT}
For \(n\) qubits:
\begin{itemize}
\item \(n\) Hadamard gates,
\item \(\frac{n(n-1)}{2}\) controlled phase gates,
\item about \(\lfloor n/2 \rfloor\) SWAP gates.
\end{itemize}

Therefore the total gate count is
\[
O(n^2).
\]

This is exponentially more efficient than applying the classical DFT matrix directly to a \(2^n\)-dimensional vector.
\end{frame}

%------------------------------------------------
\begin{frame}{Approximate QFT}
Very small-angle phase gates contribute little.

Hence we may truncate the circuit by dropping controlled \(R_k\) gates for large \(k\).

This gives the approximate QFT:
\begin{itemize}
\item fewer gates,
\item shallower circuit,
\item often still sufficient for practical algorithms.
\end{itemize}

With suitable truncation, the gate complexity can be reduced to approximately
\[
O(n \log n).
\]
\end{frame}

%------------------------------------------------
\begin{frame}{Connection to Quantum Phase Estimation I}
Suppose
\[
U \ket{\psi} = e^{2\pi i \phi}\ket{\psi}
\]
for some eigenstate \(\ket{\psi}\).

Quantum phase estimation aims to determine the phase \(\phi\).

The algorithm uses:
\begin{itemize}
\item a control register prepared in superposition,
\item controlled powers \(U^{2^j}\),
\item an inverse QFT on the control register.
\end{itemize}

The inverse QFT converts the accumulated phase information into binary digits of \(\phi\).
\end{frame}

%------------------------------------------------
\begin{frame}{Connection to Quantum Phase Estimation II}
Schematically,
\[
\frac{1}{2^{n/2}}\sum_{k=0}^{2^n-1} e^{2\pi i k\phi}\ket{k}
\]
is transformed by the inverse QFT into a computational basis state approximating the binary expansion of \(\phi\).

Thus the QFT acts as a phase decoder.

This is why it is central to:
\begin{itemize}
\item Shor's algorithm,
\item Hamiltonian eigenvalue estimation,
\item quantum simulation,
\item metrological phase extraction.
\end{itemize}
\end{frame}

%------------------------------------------------
\begin{frame}{Implementation Logic in a Notebook or Simulator}
A typical notebook implementation of the QFT proceeds as follows:

\begin{enumerate}
\item Choose the number of qubits \(n\)
\item Loop over qubits \(j=0,\dots,n-1\)
\item Apply \(H\) on qubit \(j\)
\item For each later qubit \(k>j\), apply a controlled phase \(R_{k-j+1}\)
\item At the end, reverse the qubit order with SWAP gates
\end{enumerate}

This is exactly the constructive logic derived from the product form.
\end{frame}

%------------------------------------------------
\begin{frame}{Pseudocode for QFT Construction}
\[
\begin{array}{l}
\texttt{for } j=0 \texttt{ to } n-1: \\
\qquad \texttt{apply H on qubit } j \\
\qquad \texttt{for } k=j+1 \texttt{ to } n-1: \\
\qquad\qquad \texttt{apply controlled-}R_{k-j+1}\texttt{ from } k \texttt{ to } j \\
\texttt{for } j=0 \texttt{ to } \lfloor n/2\rfloor-1: \\
\qquad \texttt{swap qubit } j \texttt{ with qubit } n-1-j
\end{array}
\]

Depending on software conventions, the control and target order may appear reversed relative to the mathematical derivation, but the logical structure is the same.
\end{frame}

%------------------------------------------------
\begin{frame}{Simulation Remarks}
When simulating the QFT numerically, it is useful to compare:
\begin{itemize}
\item the circuit output state,
\item the exact matrix formula
\[
\operatorname{QFT}_N\ket{x}
=
\frac{1}{\sqrt N}\sum_{k=0}^{N-1} e^{2\pi i xk/N}\ket{k},
\]
\item and the product-form prediction.
\end{itemize}

This provides a clean consistency check:
\begin{itemize}
\item amplitudes should have equal magnitude \(1/\sqrt N\),
\item phases should match the Fourier kernel,
\item output qubit order must be checked carefully.
\end{itemize}
\end{frame}

%------------------------------------------------
\begin{frame}{Common Source of Confusion: Bit Reversal}
A frequent issue is the ordering of bits.

The product form naturally produces the qubits in reversed order:
\[
(x_1x_2\dots x_n)\mapsto
(\text{phase involving }x_n)\otimes \cdots \otimes (\text{phase involving }x_1\dots x_n).
\]

Thus:
\begin{itemize}
\item either include final SWAP gates,
\item or accept the reversed output convention.
\end{itemize}

Many software libraries omit the final SWAPs and simply document the reversed ordering.
\end{frame}

%------------------------------------------------
\begin{frame}{Why the QFT Matters Physically}
The QFT is not just an abstract transform.

It provides a mechanism for:
\begin{itemize}
\item converting phase information into computational-basis information,
\item exploiting interference to extract periodicity,
\item connecting time evolution phases to measurable outputs.
\end{itemize}

In this sense, the QFT is a bridge between
\[
\text{dynamical phase accumulation}
\qquad\longrightarrow\qquad
\text{algorithmic readout}.
\]
\end{frame}

%------------------------------------------------
\begin{frame}{Summary}
The Quantum Fourier Transform:
\begin{itemize}
\item is the quantum analogue of the discrete Fourier transform,
\item acts as
\[
\ket{x}\mapsto \frac{1}{\sqrt N}\sum_k e^{2\pi i xk/N}\ket{k},
\]
\item factorizes into one-qubit phase states,
\item is implemented using \(H\), controlled \(R_k\), and SWAP gates,
\item has gate complexity \(O(n^2)\),
\item is central to phase estimation and many other algorithms.
\end{itemize}
\end{frame}

%------------------------------------------------
\begin{frame}{Suggested Extensions}
Natural next topics:
\begin{itemize}
\item inverse QFT in detail,
\item quantum phase estimation,
\item Shor's algorithm,
\item approximate QFT and fault-tolerant implementations,
\item explicit simulation in Qiskit or PennyLane.
\end{itemize}
\end{frame}

\end{document}
